\section{Induction Heads：实现 In-Context Nearest Neighbor 的电路}

现在回到 phase-change mystery。一个 induction head 处理模式
\[
[A][B]\ \cdots\ [A]\quad\longrightarrow\quad[B].
\]
它在当前 token 为 $A$ 时，寻找此前出现的 $A$，再读取那个旧 $A$ 后面的 $B$，提升 $B$ 的 next-token probability。

\subsection{Induction pattern}

机制通常需要两步 composition。Previous-token head 将位置 $j-1$ 的 token identity 写入位置 $j$；下一层 induction head 在当前位置的 $A$ 与候选位置所携带的 previous-token feature 之间做 QK match，并通过 copying OV 输出候选位置的当前 token $B$。

\lecturefigure{slide-49-induction-pattern.jpg}{Induction pattern：匹配此前相同 prefix，并复制它后面的 token。}{Stanford Online 官方视频 00:46:08--00:47:01。}

\begin{lstlisting}[caption={Induction-head algorithm 的概念性伪代码。}]
def induction_prediction(tokens, current_index):
    current_token = tokens[current_index]
    candidates = []
    for source in range(1, current_index):
        if tokens[source - 1] == current_token:
            candidates.append(tokens[source])
    return learned_weighted_copy(candidates)
\end{lstlisting}

\begin{importantbox}{In-context nearest neighbor}
Induction head 把 context 当作临时数据库：query 是当前 pattern，keys 是历史 pattern，value 是历史 pattern 之后发生的 token。它与 k-NN 类似，但相似度、representation、组合路径与 output transformation 都由网络学习。
\end{importantbox}

\teachervoice{Olah 的一句压缩定义是：寻找 previous copy，向前看一位，然后预测上次发生的事情会再次发生。实际 heads 可以匹配短语、token classes 或语义近邻，不只匹配字面完全相同的 $A$。}

\subsection{为什么它是 meta-learning workhorse}

简单 copying head 只会在很多可能位置盲目提升重复 token；induction circuit 则先识别“哪个过去情境与当前相似”，再复制该情境的 continuation。这个 conditional retrieval 更像一个可泛化的 context algorithm，因此能支持重复文本、格式模仿、局部函数学习与翻译对应。

\lecturefigure{slide-50-induction-workhorse.jpg}{Induction heads：meta-learning 的 workhorse，组合 pattern matching 与 copying。}{Stanford Online 官方视频 00:47:02--00:47:05。}

\lecturefigure{slide-51-multiple-induction-heads.jpg}{模型中存在多个 induction heads，覆盖不同 pattern 与 logit effects。}{Stanford Online 官方视频 00:47:06--00:47:29。}

\subsection{从 attention pattern 到 matrix signature}

对经典 induction head，OV 应具有 copying tendency；QK 应在 first-layer shifted feature 上偏好 same-token/same-pattern match。因此可以用 QK/OV traces 或 eigenvalue summaries把 heads 放入二维图：右上角同时具有 strong QK matching 与 OV copying 的 heads 是 induction candidates。

\lecturefigure{slide-52-head-taxonomy.jpg}{Induction heads、其他 heads 与 induction-ish heads 的谱/模式分类。}{Stanford Online 官方视频 00:47:30--00:47:47。}

\lecturefigure{slide-53-qk-ov-trace-scatter.jpg}{QK/OV trace scatter：用两个 circuit signatures 定位 induction candidates。}{Stanford Online 官方视频 00:47:48--00:48:15。}

\begin{warningbox}{Signature 不是定义本身}
一个 head 落在 scatter 的特定区域，只说明矩阵具有候选性质。真正 induction mechanism 还需 attention pattern、path composition、examples 与 ablation；现代模型可通过不同 basis、distributed features 或多头协作实现类似行为。
\end{warningbox}

\lecturefigure{slide-54-induction-meta-learning-summary.jpg}{Induction heads 对 meta-learning 的主要贡献与自动检测思路。}{Stanford Online 官方视频 00:48:16--00:48:21。}

\lecturefigure{slide-55-other-local-heads.jpg}{其他 heads 往往更 local，执行 token lookup 或 position tracking。}{Stanford Online 官方视频 00:48:22--00:48:33。}

\lecturefigure{slide-56-lookup-similar-context.jpg}{Induction head 的核心：查询相似历史情境中接下来发生了什么。}{Stanford Online 官方视频 00:48:34--00:48:41。}

\subsection{回到 phase-change hypothesis}

结构可行性已经具备：两层 composition 能实现 conditional lookup，induction candidates 也能从 QK/OV 与 attention pattern 中识别。于是研究假设变得具体：ICL phase change 是模型在训练中发现 induction-head circuits 的时刻。

\lecturefigure{slide-57-phase-change-hypothesis.jpg}{核心假设：induction-head 形成驱动 ICL phase change。}{Stanford Online 官方视频 00:48:42--00:49:53。}

\teachervoice{讲者再次使用 hypothesis 一词，并立即说“let's check it”。这体现了机制解释的正确节奏：先由结构提出可证伪假设，再寻找干预与跨模型证据。}

\subsection{本章小结}

Induction circuit 将 previous-token information 与下一层 copying head 组合成 $[A][B]\ldots[A]\to[B]$ 算法。它是可学习的 context nearest-neighbor mechanism，比无条件 copying 更精确，也给 phase change 提供了结构上可信的候选解释。

